
We introduced an uncertainty-aware TCN framework that treats predictive uncertainty as an internal architectural signal rather than a post-hoc add-on, accompanied by two distribution-free theoretical guarantees that anchor the operational pipeline. The framework is motivated by three sensor-specific challenges, heterogeneous reliability drift across sensor types, context-dependent anomaly scoring requirements, and the need for actionable per-sensor diagnostic output---and integrates two uncertainty-driven components: (1) the Adaptive Uncertainty-Aware Attention (AUAA) mechanism gates temporal attention by per-sensor predictive uncertainty from MC dropout; and (2) the Dynamic Weight Adapter learns context-sensitive blending of reconstruction error and uncertainty. The Chebyshev-type bound (\cref{thm:chebyshev_fp}) furnishes a distribution-free false-positive guarantee on the hybrid score, and the Monte-Carlo convergence result (\cref{thm:mc_convergence}) sets the sample-size budget $M$ for the inference pass. The architecture also includes a per-sensor attribution head, but we audit it rather than claim it as a result.

Comprehensive experiments on four-month IAQ sensor data demonstrated a $25.2\%$ MSE reduction on the held-out test split (Base $2.20\times 10^{-4}$ vs.\ Full Enhanced $1.65\times 10^{-4}$; $R^2$ 0.9988, SMAPE 11.07\%). The ablation isolates the source of this gain: the auxiliary objectives (dynamic-weight regularizer and feature-attribution loss) deliver the dominant reduction, while AUAA contributes uncertainty-gated attention at scoring time without further reducing forecasting MSE. A cross-dataset validation on the Skoltech Anomaly Benchmark (\cref{sec:skab_cross_dataset}) supplies independent, third-party ground-truth labels and confirms that the Full Enhanced configuration recovers a $7.7\times$ higher F1 ($0.477$ vs.\ $0.062$) and $12.6\times$ higher recall ($0.418$ vs.\ $0.033$) than the Base TCN on a domain (industrial water-circulation sensors) disjoint from the IAQ training distribution. The architecture additionally produces a per-sensor attribution output, although a controlled-perturbation audit (Section~4.7) shows this head is not yet causally faithful---we report it transparently as a limitation and future-work target rather than a validated explanation. The Chebyshev bound of \cref{thm:chebyshev_fp} guarantees $P(\text{FP}) \leq 1/k^2$ without distributional assumptions and is empirically respected on the benchmark. With sub-millisecond per-window inference ($0.16$\,ms at $M=20$, including the uncertainty pass) and $\sim 15\%$ parameter overhead, the framework is suitable for real-time deployment in smart buildings, environmental monitoring, and industrial process control.

\vspace{6pt}
\noindent\textbf{Operational Guarantee:} Applying the Chebyshev inequality to anomaly score distributions yields a distribution-free false-positive-rate bound: for any score distribution with finite variance and threshold \(\tau = \bar{s} + k \cdot \sigma_s\), the detector satisfies \(P(\text{FP}) \leq 1/k^2\). This enables principled threshold selection without parametric distributional assumptions, providing an operational guarantee that complements the empirical percentile-based approach.